%
% reversi_250_assemblaggio_feature.tex
%
% This file is part of the reversi program
% http://github.com/rcrr/reversi-v2
%
% Author Roberto Corradini mailto:rob_corradini@yahoo.it
% Copyright 2026 Roberto Corradini. All rights reserved.
%
% License
%
% This program is free software; you can redistribute it and/or modify it
% under the terms of the GNU General Public License as published by the
% Free Software Foundation; either version 3, or (at your option) any
% later version.
%
% This program is distributed in the hope that it will be useful,
% but WITHOUT ANY WARRANTY; without even the implied warranty of
% MERCHANTABILITY or FITNESS FOR A PARTICULAR PURPOSE. See the
% GNU General Public License for more details.
%
% You should have received a copy of the GNU General Public License
% along with this program; if not, write to the Free Software
% Foundation, Inc., 59 Temple Place - Suite 330, Boston, MA  02111-1307, USA
% or visit the site <http://www.gnu.org/licenses/>.
%

%-----------------------------------------------------------------------------
% Capitolo: Assemblaggio delle feature.
%
% Dal "pattern dentro il pattern" (la sonda delle configurazioni) ai tre
% livelli di aggregazione del modello migliore: l'ordinamento per
% prossimita' all'angolo.
% Settima fase: 2X5COR sommato agli altri 10 pattern classici
% (0056-0065): il primo livello di aggregazione, ordinato per
% prossimita' all'angolo.
% Ottava fase: il secondo livello di aggregazione (0066-0074, la base
% 2X5COR+XEDGE + un terzo pattern): l'ordinamento per prossimita'
% all'angolo.
% Nona fase: il terzo livello di aggregazione (0075-0082, la base
% 2X5COR+XEDGE+CORNER + un quarto pattern): la prossimita' ordina i
% gruppi, il campione sale a 7.82.

\chapter{Assemblaggio delle feature}
\section{Un pattern dentro l'altro: l'analisi delle configurazioni}

Prima di ogni training, le frequenze delle configurazioni sono di per
sé uno strumento di analisi. Ogni pattern possiede uno spazio di
codici — la disposizione dei dischi sulle sue celle, vuoti, neri o
bianchi, $3^k$ per $k$ celle — e il corpus ne rivela la
distribuzione. Da questa distribuzione si misurano due cose,
\emph{senza addestrare nulla}: l'\emph{espressività} del pattern, la
frazione di varianza della posizione che il pattern da solo spiega
(gli indici $\eta^2$ e $\omega^2$); e la \emph{sovrapposizione} tra
due pattern, quanto la distribuzione dell'uno determina quella
dell'altro (l'incertezza simmetrica, SU: 0 è l'ortogonalità perfetta,
1 la determinazione completa).

Il modello 0055 è il caso di prova: l'intercetta più i due pattern
EDGE e XEDGE, con la selezione REGAB su batch di dati grandi
(training sul batch 3, 1\,999\,179 righe; validazione sul batch 6).
La geometria lo decide in anticipo: le otto celle di EDGE sono otto
delle dieci celle di XEDGE. EDGE è \emph{incluso} in XEDGE.

L'analisi sulle frequenze, calcolata su un corpus di dieci milioni di
posizioni generate, conferma la geometria con i numeri.
L'espressività cresce con le celle: $\eta^2$ 0.215 per EDGE, 0.337
per XEDGE. La sovrapposizione è il dato forte: EDGE e XEDGE danno
SU 0.879 sulla stessa istanza della board — una correlazione molto
forte, l'informazione di EDGE è quasi interamente contenuta in
XEDGE — e scendono a 0.10--0.12 tra istanze diverse, dove i due
pattern sono di fatto indipendenti. La ridondanza non è un'ipotesi:
è misurata.

Il training conferma il fenomeno. I pesi del modello sono
27\,049: l'intercetta, 3\,002 configurazioni di EDGE e 24\,046 di
XEDGE, selezionate a frequenza non nulla nel corpus di
addestramento. La RMSE di validazione è 12.90 — contro il 12.91 di
XEDGE da solo: un guadagno di un centesimo, nel rumore. I 3\,002 pesi
di EDGE non aggiungono informazione: tutto quello che EDGE vede,
XEDGE lo vede già, e l'ottimizzatore lo sa. Il modello combinato è,
a tutti gli effetti, XEDGE.

Il risultato vale come metodo: prima di sommare pattern, l'analisi
delle configurazioni dice se c'è informazione nuova da guadagnare. In
questo caso lo diceva la geometria; in generale lo dicono i numeri, e
lo dicono \emph{prima} di spendere il training.

\section{Il primo livello di aggregazione: la prossimità all'angolo}

Sommare pattern è la domanda successiva: il campione della campagna,
2X5COR, portato a 9.67 con dieci volte il dato, cosa guadagna a
essere sommato agli altri pattern classici? Dieci modelli
(0056--0065) lo misurano: l'intercetta più 2X5COR più un secondo
pattern — XEDGE, CORNER, le diagonali DIAG4--DIAG8, le file
R2--R4 — sempre sulla selezione REGAB a batch grandi e con le
analisi sulle configurazioni attive.

La tabella~\ref{tab:feature-base-pattern-aggregazione} riassume i
risultati. La prima evidenza: \emph{tutte e dieci le coppie battono
2X5COR da solo} — nessun pattern lo peggiora — ma i
guadagni coprono un ordine di grandezza, da $-1.07$ a $-0.05$.

\begin{table}[ht]
\centering
\begin{tabular}{lrrrr}
\toprule
pattern & pesi $W$ & RMSE validazione & $\Delta$ su 2X5COR & SU (coppia) \\
\midrule
XEDGE  & 61\,894 & 8.60 & $-1.07$ & 0.589 \\
CORNER & 43\,004 & 9.21 & $-0.46$ & 0.601 \\
R3     & 41\,152 & 9.25 & $-0.42$ & 0.157 \\
DIAG7  & 38\,604 & 9.31 & $-0.36$ & 0.076 \\
R4     & 39\,333 & 9.32 & $-0.35$ & 0.106 \\
DIAG8  & 39\,296 & 9.33 & $-0.34$ & 0.083 \\
DIAG6  & 38\,226 & 9.41 & $-0.26$ & 0.143 \\
DIAG5  & 37\,983 & 9.49 & $-0.18$ & 0.250 \\
R2     & 41\,169 & 9.61 & $-0.06$ & 0.537 \\
DIAG4  & 37\,893 & 9.62 & $-0.05$ & 0.272 \\
\bottomrule
\end{tabular}
\caption{2X5COR sommato a ciascuno degli altri pattern classici (i
modelli 0056--0065, selezione REGAB 3/6): i pesi, la RMSE di
validazione, il guadagno su 2X5COR da solo (9.67) e la sovrapposizione
SU della coppia.}
\label{tab:feature-base-pattern-aggregazione}
\end{table}

La seconda evidenza è l'ordinamento per \emph{prossimità
all'angolo}: i guadagni decrescono al crescere della distanza del
pattern dalla regione angolare. XEDGE, che giunge a entrambi gli
angoli con la prima fila e i due contorni $B_2$ e $G_2$, apporta
$-1.07$; CORNER, il blocco $3\times3$ attorno all'angolo, $-0.46$;
la file R3, confinante con il blocco, $-0.42$; le diagonali che
nascono sul contorno dell'angolo opposto (DIAG7, DIAG8) $-0.36$ e
$-0.34$; si scende poi con la file R4 ($-0.35$) e le diagonali che
si allontanano (DIAG6 $-0.26$, DIAG5 $-0.18$), fino alla diagonale
corta e centrale DIAG4 ($-0.05$). Un caso non segue l'ordinamento e
si tiene come dato, non come legge: R2, la seconda file, confinante
anch'essa con la regione angolare, apporta $-0.06$.

La sovrapposizione SU, da sola, non ordina i guadagni: R2 ha SU
0.537, quasi come XEDGE (0.589), eppure apporta $-0.06$; DIAG7 ha
SU 0.076, il più basso della tabella, e apporta $-0.36$. Il SU
misura quanto le due distribuzioni si determinano a vicenda, in
blocco: è una misura della sovrapposizione, non dell'apporto.

Il campione aggiornato è 2X5COR+XEDGE: 8.60 di RMSE, $-67.4\%$ sul
baseline 26.42. I pesi salgono a 38\,000--62\,000 e i gap di
generalizzazione restano $+0.21$--$+0.28$: i due milioni di righe
di addestramento tengono la memorizzazione sotto controllo anche a
questa scala. È la prossimità, non la sovrapposizione, a ordinare
gli apporti: da qui partirà il secondo livello di aggregazione.

\section{Il secondo livello di aggregazione: la prossimità
all'angolo}

Completato il primo livello di aggregazione — 2X5COR più XEDGE, il
modello a 8.60 — la domanda è se i pattern restanti apportino ancora
informazione. Nove modelli (0066--0074) lo misurano: la base
2X5COR+XEDGE più un terzo pattern — CORNER, le diagonali
DIAG4--DIAG8, le file R2--R4 — sempre sulla selezione REGAB a batch
grandi; le analisi sulle configurazioni sono disattivate, a questo
punto non servono.

La tabella~\ref{tab:feature-base-pattern-secondo-livello} riassume
i risultati. La prima evidenza: tutti e nove i modelli migliorano la
base — nessun pattern la peggiora — ma gli apporti si riducono di un
ordine di grandezza rispetto al primo livello (da $-1.07$ a
$-0.02$--$-0.49$): la massa predittiva della posizione è concentrata
nella regione angolare, e il primo livello l'ha in larga parte
raccolta.

\begin{table}[ht]
\centering
\begin{tabular}{lrrrr}
\toprule
pattern & pesi $W$ & RMSE validazione & $\Delta$ sulla base & gap (val $-$ train) \\
\midrule
CORNER & 67\,050 & 8.11 & $-0.49$ & $+0.29$ \\
R3     & 65\,198 & 8.15 & $-0.45$ & $+0.29$ \\
R4     & 63\,379 & 8.22 & $-0.38$ & $+0.28$ \\
DIAG7  & 62\,650 & 8.22 & $-0.38$ & $+0.29$ \\
DIAG8  & 63\,342 & 8.24 & $-0.36$ & $+0.27$ \\
DIAG6  & 62\,272 & 8.32 & $-0.28$ & $+0.29$ \\
DIAG5  & 62\,029 & 8.41 & $-0.19$ & $+0.28$ \\
DIAG4  & 61\,939 & 8.56 & $-0.04$ & $+0.29$ \\
R2     & 65\,215 & 8.58 & $-0.02$ & $+0.31$ \\
\bottomrule
\end{tabular}
\caption{La base 2X5COR+XEDGE (8.60) più un terzo pattern (i modelli
0066--0074, selezione REGAB 3/6): i pesi, la RMSE di validazione,
l'apporto sulla base e il gap di generalizzazione, in ordine
decrescente di apporto.}
\label{tab:feature-base-pattern-secondo-livello}
\end{table}

La seconda evidenza è l'ordinamento per \emph{prossimità
all'angolo}. Gli apporti decrescono al crescere della distanza del
pattern dalla regione angolare: CORNER, il blocco $3\times3$
attorno all'angolo, $-0.49$; la file R3, confinante con il blocco,
$-0.45$; la file R4, il gradino successivo, $-0.38$; le diagonali
che nascono sul contorno dell'angolo opposto (DIAG7, DIAG8) $-0.38$
e $-0.36$; le diagonali che se ne allontanano (DIAG6, DIAG5)
$-0.28$ e $-0.19$; la diagonale corta e centrale (DIAG4) $-0.04$. È
lo stesso principio che governa la campagna fin dalle sue origini,
ora formulato in termini di prossimità: a parità di tutto il resto,
l'apporto è maggiore per il pattern che insiste più direttamente
sulla regione angolare.

Un caso non segue l'ordinamento e si tiene come dato, non come
legge: R2, la seconda file, confinante anch'essa con la regione
angolare, apporta $-0.02$.

Il campione aggiornato è 2X5COR+XEDGE+CORNER: 8.11 di RMSE,
$-69.3\%$ sul baseline 26.42, 67\,050 pesi, gap di
generalizzazione $+0.29$: la memorizzazione resta contenuta anche a
questa scala, grazie ai due milioni di righe di addestramento. Il
modello migliore della campagna è costruito, per la sua parte di
pattern, quasi interamente sulla regione angolare.

\section{Il terzo livello di aggregazione: la prossimità
all'angolo}

Completato il secondo livello — la base 2X5COR+XEDGE+CORNER, il
modello a 8.11 — otto modelli (0075--0082) misurano il prossimo
gradino: la base più un quarto pattern — le diagonali
DIAG4--DIAG8, le file R2--R4 — sempre sulla selezione REGAB a batch
grandi; le analisi sulle configurazioni restano disattivate.

La tabella~\ref{tab:feature-base-pattern-terzo-livello} riassume i
risultati. Sette pattern apportano un miglioramento, uno non
aggiunge nulla: la diagonale corta e centrale DIAG4 resta
esattamente sulla base. Gli apporti si riducono ancora rispetto al
secondo livello (da $-0.49$ a $-0.29$): la massa predittiva della
posizione resta concentrata nella regione angolare, e ogni livello
ne estrae la frazione residua.

\begin{table}[ht]
\centering
\begin{tabular}{lrrrr}
\toprule
pattern & pesi $W$ & RMSE validazione & $\Delta$ sulla base & gap (val $-$ train) \\
\midrule
R4     & 68\,535 & 7.82 & $-0.29$ & $+0.28$ \\
DIAG7  & 67\,806 & 7.84 & $-0.27$ & $+0.29$ \\
R3     & 70\,354 & 7.88 & $-0.23$ & $+0.29$ \\
DIAG6  & 67\,428 & 7.91 & $-0.20$ & $+0.28$ \\
DIAG8  & 68\,498 & 7.97 & $-0.14$ & $+0.29$ \\
DIAG5  & 67\,185 & 8.05 & $-0.06$ & $+0.29$ \\
R2     & 70\,371 & 8.09 & $-0.02$ & $+0.30$ \\
DIAG4  & 67\,095 & 8.11 & $0.00$  & $+0.30$ \\
\bottomrule
\end{tabular}
\caption{La base 2X5COR+XEDGE+CORNER (8.11) più un quarto pattern
(i modelli 0075--0082, selezione REGAB 3/6): i pesi, la RMSE di
validazione, l'apporto sulla base e il gap di generalizzazione, in
ordine decrescente di apporto.}
\label{tab:feature-base-pattern-terzo-livello}
\end{table}

La prossimità all'angolo continua a ordinare i \emph{gruppi}. Il
gruppo di testa, formato da quattro pattern adiacenti alla regione
angolare — le file R4 e R3, vicine al blocco A1, e le diagonali
nate sul bordo, dal lato di H1 (DIAG7, DIAG6) — copre la fascia
$-0.29$--$-0.20$; il fondo è occupato dai pattern centrali (DIAG5,
R2, DIAG4), tra $-0.06$ e $0.00$. Dentro il gruppo, l'ordine fine
non è rigido: a questo livello R4 supera R3, mentre al secondo
livello era R3 a comandare. DIAG8 si colloca a metà della scala
($-0.14$): la diagonale nasce nell'angolo H1, che XEDGE ha già
acquistato al primo livello, e il suo apporto residuo è la coda
verso l'angolo lontano, che decresce col distacco dalla regione già
coperta.

Il campione aggiornato è 2X5COR+XEDGE+CORNER+R4: 7.82 di RMSE,
$-70.4\%$ sul baseline 26.42, 68\,535 pesi, gap di
generalizzazione $+0.28$: la memorizzazione resta contenuta anche a
questa scala. I quattro pattern del modello migliore della campagna
sono tutti, ciascuno a modo suo, sulla regione angolare.
